\begin{comment}

\subsection{What's Beyond?}
\label{sec:whats-beyond}

\dSe{
    Hello.
}

\dMax{
    \ldots{} hey.
}

\dSe{
    Now that we're here, tell me what you see.
    What's beyond?
}

\dMax{
    \ldots{}
}

\dMax{
    I'm not really sure.
    It's\ldots{} overwhelming to think about.
}

\dSe{
    Do you know why?
}

\dMax{
    Maybe.
    But I'm reticent to say.
}

\dMax{
    It would violate my sense of humility.
}

\dSe{
    And you mean that\ldots{}
}

\dMax{
    Both ways.
}

\dSe{
    So you don't want to discuss it then?
}

\dMax{
    \ldots{} I didn't say that.
}

\dSe{
    What's stopping you?
    And no evasion this time.
}

\dMax{
    \ldots{}
    I'm scared of what will happen if I'm right.
    That I'll lose the ladder I built getting this far.
}
\dMax{
    In some ways I already feel it happening.
}

\dSe{
    Then hold on tighter. You know how, after all.
}
\dSe{
    But, that first thing you said.
    If you're right, then this is big, sure.
    Why will that compromise you?
}

\dMax{
    Because the only reason I've gotten here is knowing that I can be wrong, that I'm probably wrong.
    Not just knowing\ldots{}
    More like living.
}
\dMax{
    If I'm right, then everything changes and it get a lot harder.
}

\dSe{
    What gets harder?
}

\dMax{
    Maintaining my humility.
    Even talking about it feels like it's at risk.
    Like, who am I to think I know anything about humility?
}

\dMax{
    It's like some fragile thing that I can't even bear to look at, lest somehow my gaze alone is enough to shatter it.
}

\dMax{
    How can I be honest about what I actually think without contradicting the very values that made UT possible?
}

\dMax{
    It feels like a catch-22.
}

\dSe{
    Stop.
    What do we do when there's a contradiction?
}

\dMax{
    Right, check my premises.
}

\dSe{
    \ldots{}
}

\dMax{
    Well, being honest about what I think is more important than being right about humility.
    Or more foundational, or whatever.
}

\dSe{
    Mm. And what is your sense of humility?
}

\dMax{
    Knowing that I'm probably wrong, and that I'm not as good a thinker as I think I am.
    But always seeking and desiring criticism, for it's a gift.
    It's not enough just to know this in a superficial level, but I have to honestly live it, too.
    So some part of it is temperance.
}

\dSe{
    And you've had this idea for how long?
}

\dMax{
    A few years.
}

\dSe{
    You know what definitely will make maintaining your sense of humility harder?
}

\dSe{
    Holding on to it if it's wrong.
}

\dMax{
    Why do you think it's wrong?
}

\dSe{
    Do you really expect me to do all the work?
}

\dMax{
    It's wrong because I hold on to it like a life-raft.
    Which is another embodyment of believing in perfection and infalliblity.
}

\dSe{
    ``Because''?
}

\dMax{
    Right, it's wrong \emph{to} hold on to it like a life-raft.
    We should let ideas die in our place.
}

\dSe{
    xxx
}

\hrulefill

\dSe{xxx}

\hrule



% ---

\dMax{
    I'm somewhat worried what people will think, too.
    You know, people I care about.
}

\dSe{
    Because it's a blockchain thing.
}

\dMax{
    Yeah.
    I can't really blame people for being negative about it.
    I mean, look at the state of things.
}

\dSe{
    And what have you done to deserve the same judgement?
}

\dSe{
    Did you stick to your principles?
}

\dMax{
    Yes.
}

\dSe{
    Walk me through the implications of that answer.
}

\dMax{
    Well, maybe I think so little of my friends that I don't trust them to stick to their principles.
    Or, maybe I'm secondhanded and will live by their values instead of mine.
    But I don't want to live like that.
}

\dSe{
    Then don't. Live for yourself.
    Stop being worried about all the ways you might fail.
    If you do, stick to your principles and you'll see it.
    Then you can fix it.
    Stop worrying about your friends, they can take care of themselves.
    \\
    Besides, how can living for yourself scare you when you know the alternative?
}

\dMax{
    Because it's the first time I'll really have to.
}

\dSe{
    \ldots
    So, what do we do with hard things?
    Go towards them, or away?
}

\dMax{
    Towards.
}

\end{comment}
